% !TeX root = ../教材题库_全部条目_转写工作稿.tex
% Source: XJ-B1-C06-HW-P222-Y0836
% 第6章 / 6.2 抽样 / 教材 p.222 / PDF 第229页
% OCR未校对：公式与文字均可能有误；未完成数学TeX转写；禁止直接发给学生
% 原题及相邻原解截图（核对用）：../题目原图/XJ-B1-C06-HW-P222-Y0836-01.webp
% 本文件只在转写完成后作为题干使用；OCR 条目尚不是可靠题干。
\begin{bankitem}{习题6.2 第7题}{第6章 / 6.2 抽样 / 教材 p.222 / PDF 第229页}{XJ-B1-C06-HW-P222-Y0836}{OCR未校对：公式与文字均可能有误；未完成数学TeX转写；禁止直接发给学生}
\begin{ocrdraft}
[教材p.222 / PDF第229页]
\par
7．一家科研单位最近组织了一次健身活动，活动分为登山组和游泳组，且每
\par
个职工至多参加其中一组．在参加活动的职工中，青年人占42，5％，中年人占
\par
37．5％，老年人占20％．登山组的职工占参加活动总人数的三分之一，且该组中，
\par
青年人占50％，中年人占30％，老年人占20％．为了解各组不同年龄层次的职工
\par
对本次活动的整体满意程度，现用分层抽样的方法从参加活动的全体职工中抽取一
\par
个容量为200的样本．试确定：
\par
（1）游泳组中，青年人、中年人、老年人分别所占的比例；
\par
（2）游泳组中，青年人、中年人、老年人分别应抽取的人数．
\end{ocrdraft}
\end{bankitem}
